\begin{abstract}
We study a randomized trial in which treatment and outcome remain linked to a deterministic label of the assignment probability. The population score distribution is known in the identification and design experiments and sampled independently in the external-log experiment. Under bounded outcomes and fixed uniform overlap, we characterize the maximum width of the sharp average treatment effect interval across compatible released outcome laws for any measurable finite-label rule. Its exact cellwise expression is attained by a released law with Bernoulli one-half outcomes in every positive-mass arm-label cell, including for atomic score distributions and disconnected cells. Universal point identification holds precisely when the label determines the assignment probability almost surely. With a budget of \(K\) labels, optimal worst-case ambiguity has order \(K^{-1}\) over score laws with a common positive density lower bound; an equal-width construction gives the upper bound for every score law. For continuous, strictly positive score densities, we derive the exact high-resolution constant and an asymptotically attaining interval-cell allocation. For trial size \(n\) and a known score law, the minimax expected length of a uniformly honest interval has order \(K^{-1}+n^{-1/2}\) over the lower-density class. With a fixed release and an independent score log of size \(m\), a projection interval attains excess-length order \(n^{-1/2}+m^{-1/2}\); separate lower bounds establish both sampling directions at a fixed positive trial-to-log size ratio.
\end{abstract}

\section{Introduction}

Randomized assignment can use a probability that varies across people even when a released trial file records only treatment and outcome. We consider a release that attaches to each trial row a deterministic label of that assignment probability. In the identification, design, and known-score inference experiments, the analyst knows the population distribution of assignment probabilities. In the external-log experiment, that distribution is learned from an independent score sample. This observation scheme retains the link among label, treatment, and outcome, while the population score distribution supplies the probabilities represented within each label. The question is how much information the labels preserve about the average treatment effect (ATE), and how to choose them when their number is limited.

Under bounded potential outcomes and a fixed uniform overlap margin, the sharp ATE interval for a specified compatible released law follows from within-cell quantile rearrangement. We use the bounded-outcome specialization of \citet[Theorem 3.2 and eq. (3.1)]{FanShermanShum2014} for this fixed-law interval; \cref{obj:lem:fixed-released-law-baseline} records the resulting endpoints. Our release-design criterion then takes the largest interval width across the released outcome laws compatible with a given score distribution and label rule. \Cref{obj:thm:all-label-ambiguity} expresses that width as a sum of paired absolute-deviation losses over label cells and constructs two causal laws attaining its endpoints under the same released law. The construction applies to arbitrary score distributions on the overlap interval and measurable cells, including atoms and disconnected sets. \Cref{obj:thm:universal-point-identification} gives the corresponding disclosure criterion: worst-case ambiguity is zero exactly when the label determines the score almost surely.

This criterion turns finite-label release into a design problem measured in ATE units. For a budget of \(K\) labels, \cref{obj:thm:k-label-rate} bounds the smallest worst-case ambiguity above and below at order \(K^{-1}\) when the score density has a common positive lower bound. Its equal-width upper construction applies to every score distribution on the overlap interval. When the density is continuous and strictly positive, \cref{obj:thm:sharp-high-resolution-constant} gives the exact leading constant. Writing \(\varepsilon\) for the overlap margin, \(e\) for a score value, and \(f\) for the score density, an asymptotically attaining sequence places interval boundaries at equal increments of \(\int_\varepsilon^e \sqrt{f(t)/[t(1-t)]}\,dt\). Thus label allocation responds to both score mass and the sensitivity of inverse assignment weights.

The design criterion also governs inference from independent released trial rows. With a known score law and a jointly chosen release and uniformly honest interval, \cref{obj:thm:minimax-honest-length} establishes minimax expected length of order \(K^{-1}+n^{-1/2}\) over the common positive lower-density class, where \(n\) is trial size. An equal-width release and midpoint inverse-probability weights attain the upper order. A label budget proportional to \(\sqrt n\) balances ambiguity from score variation within cells with trial sampling variation. This connects the partial-identification criterion to confidence-set analysis \citep{ImbensManski2004,Stoye2009,ChernozhukovHongTamer2007,KaidoMolinariStoye2019}.

We also study a fixed finite-label release when the score distribution is learned from an independent log of \(m\) scores. The criterion here is worst-case expected positive excess length beyond the population sharp interval. For every measurable fixed release and score law on the overlap interval, \cref{obj:thm:external-score-root-order} constructs a uniformly honest projection interval with excess length bounded at order \(n^{-1/2}+m^{-1/2}\). A trial-sample lower bound holds at every positive pair of sample sizes, and a positive score-log lower bound holds asymptotically uniformly over trial sizes. Together they establish that order when the trial-to-log size ratio tends to a positive constant. The result includes atomic score laws and quantile functions with flat pieces.

The setting connects partial identification under combined marginals \citep{Manski1990,Tamer2010,Torgovitsky2019} with statistical matching \citep{DOrazio2006,Conti2017} and propensity-score subclassification \citep{Cochran1968,RosenbaumRubin1984}. The released label specifies a deterministic score cell, making within-cell association and the choice of cells central to the analysis. We document the scope of the Lean 4 verification in the appendix. We place the observation scheme among related work in the next section. The paper then develops the trial setup and sharp interval, characterizes information retained by labels, studies finite-label design and honest intervals, analyzes the independent score-log experiment, and closes with discussion and open questions.
